\chapter{Outlook}
In this work, the results for raw transverse momentum spectra for charged cascade particles $\Xi^{\pm}$ and $\Omega^{\pm}$ in proton-lead collisions at
$\sqrt{s_{\textnormal{\sc nn}}}$ = 8.16 TeV in ALICE were presented.

Developing a basic theoretical understanding supported by experimental results pertaining to Quark-Gluon Plasma and its signatures were described. A brief
overview of the experimental setup and the detector assemblies relevant to the analysis was also detailed. This was followed by details of the analysis,
including data collection and reconstruction, data sample selection, event and candidate selection criteria, and an explanation of the signal extraction
method employed in this analysis. The optimised selection criteria was tabulated and raw (uncorrected) transverse momentum spectra for the cascade particles
distributed over different multiplicity classes and $p_{\textnormal{\sc t}}$ ranges was shown.

Future aspects of this work include pile-up estimation and correction using all available data samples in this dataset, detector efficiency estimation and
acceptance corrections using the Monte Carlo productions and obtaining the efficiency-corrected transverse momentum spectra. This spectra can be
compared to available $\Xi^{\pm}$
and $\Omega^{\pm}$ $p_{\textnormal{\sc t}}$ spectra in p-Pb collisions at different energies. Finally, systematic uncertainties estimation for chosen selection
criteria and error estimation will be performed.



\chapter*{Súhrn}
V tejto práci boli prezentované nekorigované spektrá priečnych hybností nabitých častíc $\Xi^{\pm}$ a $\Omega^{\pm}$ v zrážkach protónov s olovom pri 
$\sqrt{s_{\textnormal{\sc nn}}}$ = 8,16 TeV v experimente ALICE na Veľkom hadrónovom zrážači (LHC) v CERN.

Kvalitatívne boli opísané teoretické základy pre pochopenie fyziky kvarkovo-gluónovej plazmy (QGP) a uvedený stručný prehľad niektorých experimentálnych 
signatúr QGP. Bol opísaný detektorový systém experimentu ALICE so zameraním na kľúčové detektory potrebné pre rekonštrukciu rozpadov nabitých multipodivných 
hyperónov. Nasledoval podrobný opis analýzy dát -- výber vhodných zrážok, výber selekčných kritérií pre kandidátov na rozpad a popis extrakcie signálu z 
rozdelenia invariantných hmotností rozpadových produktov. Finálne selekčné kritériá boli uvedené v tabuľke a nekorigované spektrá priečnych hybností multipodivných 
hyperónov $\Xi^{\pm}$ a $\Omega^{\pm}$ boli zobrazené pre jednotlivé multiplicitné triedy.

V budúcnosti bude analýza rozšírená o korekciu príspevku z tzv. pile-up eventov (t.j. viac zrážok zrekonštruovaných ako jedna zrážka), budú pridané datasety 
s vyššou frekvenciou zrážok a teda s vyšším obsahom pile-up eventov. Pomocou zrekonštruovaných Monte Carlo dát bude odhadnutá korekcia na detektorovú akceptanciu 
a účinnosť rekonštrukcie kaskádnych rozpadov $\Xi^{\pm}$ a $\Omega^{\pm}$. Nakoniec budú odhadnuté systematické chyby z rôznych uvažovaných zdrojov.

Výsledné korigované spektrá priečnych hybností nabitých častíc $\Xi^{\pm}$ a $\Omega^{\pm}$ budú porovnané so spektrami z p-Pb zrážok avšak pri inej zrážkovej 
energii. Taktiež budú porovnané s inými zrážkovými systémami, ako napr. vysokomultiplicitné protón-protónové zrážky alebo Pb—Pb zrážky v závislosti od multiplicity 
nabitých častíc v evente.

